% =============================================================================
% main.tex —— 唯一编译入口
% -----------------------------------------------------------------------------
% 编译：tectonic main.tex      （生成 main.pdf）
%
% 加载顺序不可随意调换：
%   1) 先读 style-*.tex：它用 \PassOptionsToClass 传字号，并提供
%      \SlideGeometry / \ApplySlideStyle / \InkMargin / \SourceNoteSize，
%      其中 \SlideGeometry 必须在 \documentclass 之后立刻执行。
%   2) 再读公共设置：base → colors → layout → toc → blocks → commands。
%   3) 最后 \ApplySlideStyle 施加版式字号与页脚，再读封面元信息。
% =============================================================================

% -----------------------------------------------------------------------------
% 1. 选择版式：small（默认，16 × 9 cm）或 large（24 × 13.5 cm）
%    只改下面这一行即可切换；两版共用正文、字体、配色与图片。
% -----------------------------------------------------------------------------
% =============================================================================
% style-small.tex —— 小版版式（本模板默认）
% -----------------------------------------------------------------------------
% 职责：只定义某一版「页面几何 + 文档字号 + 边距 + 页脚」这四件事。
%       公共配色、Block、页标题、过渡页等全部位于 config/ 下的其他文件，
%       因此切换版式不会影响观感，只会改变物理尺寸与字号。
%
% 加载时机：必须在 \documentclass{beamer} 之前 \input。
%   本文件用 \PassOptionsToClass 传字号，并在类加载后由 \SlideGeometry 设定纸张。
%
% 小版特征：16 × 9 cm（面积约为大版的 44%）、11pt；
%   页脚只有页码，没有横线和额外竖直留白。
% =============================================================================

% 文档字号：11pt（作用于所有相对字号，例如 \large、block 标题）
\PassOptionsToClass{11pt}{beamer}

% 页面尺寸：在 main.tex 于 \documentclass 之后立刻执行
\newcommand{\SlideGeometry}{\geometry{paperwidth=16cm,paperheight=9cm}}

% 左右边距：标题、上横线、正文共享同一数值，改这里即可整体收放
\newcommand{\InkMargin}{0.9cm}

% \sourceNote 的字号：小版用 \scriptsize
\newcommand{\SourceNoteSize}{\scriptsize}

% 版式相关字号与页脚：在导言区末尾由 main.tex 调用，此时公共模板已就绪
\newcommand{\ApplySlideStyle}{%
  \setbeamerfont{title}{size=\LARGE}
  \setbeamerfont{frametitle}{size=\fontsize{17}{21}}
  \setbeamerfont{section title}{size=\LARGE}
  \setbeamerfont{block title}{size=\fontsize{12}{15}}
  \setbeamerfont{subsection in toc}{size=\scriptsize}
  \setbeamerfont{footnote}{size=\scriptsize}
  \setbeamerfont{page number in head/foot}{size=\scriptsize}
  % 小版页脚：只有页码，无横线、无底部留白
  \setbeamertemplate{footline}{%
    \begin{beamercolorbox}[wd=\paperwidth,leftskip=\InkMargin,rightskip=\InkMargin]{page number in head/foot}
      \InkPageNumber
    \end{beamercolorbox}}}

\documentclass[aspectratio=169,t]{beamer}
\SlideGeometry

% -----------------------------------------------------------------------------
% 2. 导言区（公共设置，与版式无关）
% -----------------------------------------------------------------------------
% =============================================================================
% base.tex —— 基础层：宏包、中文字体、表格、文献、正文用自定义命令
% -----------------------------------------------------------------------------
% 加载顺序见 main.tex；本文件在 Beamer 类之后、其余公共配置之前加载。
% 本文件只依赖「已经载入 beamer」这一前提，不依赖配色与版式设置。
% =============================================================================

% -----------------------------------------------------------------------------
% 1. 中文与字体
% -----------------------------------------------------------------------------
% Tectonic / XeLaTeX 编译；不让 fontspec 接管数学字体，避免与 amssymb 冲突。
% ctex 在 fontset=none 下只做中文排版机制，字体由下面的 fontspec 命令显式指定。
\PassOptionsToPackage{no-math}{fontspec}
\usepackage[UTF8,fontset=none]{ctex}
\defaultfontfeatures{Ligatures=TeX}

% 字体策略：优先使用随项目携带的思源宋体/黑体（fonts/ 目录，无需系统安装）；
% 如果这些 OTF 被删掉，则回退到 TeX 发行版自带的 Fandol。
% 四个文件必须全部存在，避免字体的 Regular / Bold 指向缺失文件。
% 中文没有西文意义上的 italic：斜体统一映射到 Fandol 楷体，
% 粗斜体保持黑体/宋体粗体，避免出现伪斜体。
\newif\ifInkBundledFonts
\InkBundledFontstrue
\newcommand{\InkCheckFont}[1]{\IfFileExists{fonts/#1}{}{\InkBundledFontsfalse}}
\forcsvlist{\InkCheckFont}{SourceHanSerifCN-Regular.otf,SourceHanSerifCN-Bold.otf,SourceHanSansCN-Regular.otf,SourceHanSansCN-Bold.otf}
\ifInkBundledFonts
  % 完整思源字体：衬线、无衬线与中文等宽族
  \setCJKmainfont{SourceHanSerifCN-Regular.otf}[Path=fonts/,
    BoldFont=SourceHanSerifCN-Bold.otf,
    ItalicFont=FandolKai-Regular.otf,ItalicFeatures={Path=},
    BoldItalicFont=SourceHanSerifCN-Bold.otf]%
  \setCJKsansfont{SourceHanSansCN-Regular.otf}[Path=fonts/,
    BoldFont=SourceHanSansCN-Bold.otf,
    ItalicFont=FandolKai-Regular.otf,ItalicFeatures={Path=},
    BoldItalicFont=SourceHanSansCN-Bold.otf]%
  % 等宽场景（代码、路径）用思源黑体，保证中文可读
  \setCJKmonofont{SourceHanSansCN-Regular.otf}[Path=fonts/,
    BoldFont=SourceHanSansCN-Bold.otf,
    ItalicFont=FandolKai-Regular.otf,ItalicFeatures={Path=},
    BoldItalicFont=SourceHanSansCN-Bold.otf]%
\else
  % 思源字体不完整：整套回退到 Fandol
  \setCJKmainfont{FandolSong-Regular.otf}[BoldFont=FandolSong-Bold.otf,
    ItalicFont=FandolKai-Regular.otf,BoldItalicFont=FandolSong-Bold.otf]%
  \setCJKsansfont{FandolHei-Regular.otf}[BoldFont=FandolHei-Bold.otf,
    ItalicFont=FandolKai-Regular.otf,BoldItalicFont=FandolHei-Bold.otf]%
  \setCJKmonofont{FandolFang-Regular.otf}[BoldFont=FandolHei-Bold.otf,
    ItalicFont=FandolKai-Regular.otf,BoldItalicFont=FandolHei-Bold.otf]%
\fi
% 西文保留 Latin Modern（beamer 默认），无需额外设置。

% -----------------------------------------------------------------------------
% 2. 表格
% -----------------------------------------------------------------------------
% booktabs 提供 \toprule/\midrule/\bottomrule 三线表；
% tabularx 提供自适应列宽，下面定义 Y 列 = 自动宽度 + 左对齐（避免两端对齐的稀疏感）。
\usepackage{booktabs,tabularx}
\newcolumntype{Y}{>{\raggedright\arraybackslash}X}
% 表头按单元格显式标记；普通表格、自动宽度表格与合并单元格均可复用。
\newcommand{\tablehead}[1]{\textbf{#1}}
% 表格统一使用衬线体：中文为思源宋体、西文为 Latin Modern Roman。
% 用环境分组包住，保证表外的正文与来源说明仍用原字体。
\newcommand{\InkTableStyle}{\rmfamily\renewcommand{\arraystretch}{1.2}}
\AtBeginEnvironment{tabular}{\InkTableStyle}
\AtBeginEnvironment{tabularx}{\InkTableStyle}

% -----------------------------------------------------------------------------
% 3. 图片与链接
% -----------------------------------------------------------------------------
% 图片统一放在 figures/，正文里只写文件名即可
\graphicspath{{figures/}}
% 网址与正文用同一套字形，避免 URL 突然换字体
\urlstyle{same}

% -----------------------------------------------------------------------------
% 4. 引用类文字的专用字体（Times 风格）
% -----------------------------------------------------------------------------
% 参考文献、引文块与脚注使用 Times 风格衬线体，
%       而正文西文仍保留 Latin Modern。
% 说明：TeX 发行版不自带 Times New Roman（商业字体），这里使用其度量兼容的
%       TeX Gyre Termes（Nimbus Roman No9 的再发布版），字形与 Times New Roman
%       基本一致，且随 Tectonic / TeX Live 一起分发，无需额外安装。
% 使用文件名加载，避免依赖系统字体索引。
% 若某套发行版只提供 X 变体，则回退到 TeXGyreTermesX-*.otf，保证仍可编译。
\IfFileExists{texgyretermes-regular.otf}{%
  \newfontfamily\TimesCite{texgyretermes-regular.otf}[%
    BoldFont=texgyretermes-bold.otf,
    ItalicFont=texgyretermes-italic.otf,
    BoldItalicFont=texgyretermes-bolditalic.otf]%
}{%
  \newfontfamily\TimesCite{TeXGyreTermesX-Regular.otf}[%
    BoldFont=TeXGyreTermesX-Bold.otf,
    ItalicFont=TeXGyreTermesX-Italic.otf,
    BoldItalicFont=TeXGyreTermesX-BoldItalic.otf]%
}

% 引用块环境：沿用 LaTeX 标准名 quote / quotation，但换成 Times 风格字体。
% quote 适合较短引文（整段缩进、无首行缩进）；
% quotation 适合多段落长引文（整段缩进且首行再缩进）。
% 保存原始环境命令，包裹字体切换并保持原有缩进和分段规则。
\let\origquote\quote
\let\endorigquote\endquote
\let\origquotation\quotation
\let\endorigquotation\endquotation
\renewenvironment{quote}
  {\origquote\TimesCite\small}
  {\endorigquote}
\renewenvironment{quotation}
  {\origquotation\TimesCite\small}
  {\endorigquotation}

% -----------------------------------------------------------------------------
% 5. 参考文献
% -----------------------------------------------------------------------------
% 后端固定为 BibTeX：Tectonic 内置 BibTeX，但**不提供 biber**，
% 所以这里不能改成 backend=biber（会以 "program not found" 失败）。
% 若将来需要 biber 的完整多语种/中文作者支持，必须改用带 biber 的
% XeLaTeX 工具链（见 README「参考文献」一节）。
% 作者年份制、按姓名-年份-标题排序，正文引用最多 2 位作者，文献表最多 6 位。
\usepackage[backend=bibtex,style=authoryear,sorting=nyt,
  maxcitenames=2,maxbibnames=6,giveninits=true,uniquename=init,
  doi=true,url=true,isbn=false]{biblatex}
\addbibresource{references/references.bib}

% 参考文献表字体：必须通过 biblatex 的 \bibfont 钩子设置。
% 在 \bibfont 中设置，避免条目排版时重置回正文字体。
\renewcommand{\bibfont}{\TimesCite}

% csquotes：提供语言自适应的引号命令 \enquote{...} 与长引文 \blockquote{...}，
% 免去手工写 ``...'' 这类 TeX 引号；autostyle 让它跟随文档语言选择引号形式。
\usepackage[autostyle]{csquotes}

% -----------------------------------------------------------------------------
% 6. 正文中使用的可复用命令
% -----------------------------------------------------------------------------
% 突出重点：使用主题结构色 + 粗体。
\newcommand{\keypoint}[1]{\structure{\textbf{#1}}}

% 独立段落注释：无标题栏的浅灰圆角盒，小一号灰色直立正文。
% 原生圆角盒左右各外扩 4bp，内缩后与当前栏边界对齐。
\newcommand{\slideNote}[1]{%
  \par\smallskip\begingroup
  \noindent\hspace*{4bp}%
  \begin{minipage}{\dimexpr\linewidth-8bp\relax}
    \begin{beamercolorbox}[wd=\linewidth,sep=0.65em,rounded=true,shadow=false]{annotation}
      \small\sffamily\mdseries\upshape\raggedright
      \setlength{\parskip}{0.25em}#1\par
    \end{beamercolorbox}
  \end{minipage}\par\endgroup\smallskip}

% 来源说明（图表出处、数据来源）：小字号、灰色、左对齐，无脚注编号。
% 内容通常是 \parencite 之类的引用，因此一并使用 Times 风格字体。
% 在调用处就地排版，建议写在 frame 末尾。
\newcommand{\sourceNote}[1]{%
  \par\smallskip
  {\raggedright\SourceNoteSize\TimesCite\color{AcademicGray}#1\par}}

% 插图：第一参数为高度上限（默认 0.58\textheight），宽度自动适应单栏或双栏，
% 并保持原始比例。例：\slideimage[0.45\textheight]{result.png}
\newcommand{\slideimage}[2][0.58\textheight]{%
  \includegraphics[width=\linewidth,height=#1,keepaspectratio]{#2}}
       % 宏包、中文字体、表格、文献、正文命令
% =============================================================================
% colors.tex —— 墨线主题配色
% -----------------------------------------------------------------------------
% 只有四组基础色 + 三个 Block 语义色。所有颜色集中在此，
% 换配色时不需要碰版式、Block 结构或任何正文文件。
% 语义约定（正文里按语义选色，不要记十六进制值）：
%   Ink           深墨色，标题与正文
%   AcademicBlue  主题蓝，结构色（页标题横线、编号、\alert、\keypoint、普通 Block）
%   AcademicGray  辅助灰，来源说明、副信息
%   InkLine       浅灰线，仅用于大版页脚横线与图片占位框
%   InkAlert      砖红，警示 / 限制 / 注意事项（alertblock）
%   InkExample    灰绿，算例 / 实验设置 / 应用场景（exampleblock）
% 配色思路参考 UCL Beamer 的深色标题栏 + 同色系浅色正文背景，
% 以及 Metropolis 对「警示」与「示例」的语义区分；
% 但只使用 Beamer / xcolor 原生能力，不依赖任何外部主题。
% =============================================================================

\usetheme{default}

% 基础色
\definecolor{AcademicBlue}{HTML}{345D82}   % 主题蓝：结构、强调、标题横线
\definecolor{AcademicGray}{HTML}{6E7A85}   % 辅助灰：来源说明与小字
\definecolor{Ink}{HTML}{17212B}            % 深墨色：标题与正文
\definecolor{InkLine}{HTML}{AAB4BB}        % 线条灰：页脚横线、占位框

% Block 语义色：说明用墨蓝、警示用砖红、示例用灰绿
\definecolor{InkAlert}{HTML}{9B493F}       % 警示块
\definecolor{InkExample}{HTML}{35695A}     % 示例块

% 全局文字与结构色
\setbeamercolor{normal text}{fg=Ink,bg=white}
\setbeamercolor{structure}{fg=AcademicBlue}
\setbeamercolor{title}{fg=Ink}
\setbeamercolor{section title}{fg=Ink}
\setbeamercolor{frametitle}{fg=Ink,bg=white}
% \alert 与 \keypoint 都用主题蓝加粗
\setbeamercolor{alerted text}{fg=AcademicBlue}
\setbeamerfont{alerted text}{series=\bfseries}

% Block 三色：标题栏用 88% 混色 + 白字，正文背景用 16% 混色，
% 数值越小越淡。加星号清除默认的颜色继承链，
% 避免「块内强调色」与「标题色」互相引用造成颜色串味。
\setbeamercolor*{block title}{fg=white,bg=AcademicBlue!88!white}
\setbeamercolor*{block body}{fg=Ink,bg=AcademicBlue!16!white}
\setbeamercolor*{block title alerted}{fg=white,bg=InkAlert!88!white}
\setbeamercolor*{block body alerted}{fg=Ink,bg=InkAlert!16!white}
\setbeamercolor*{block title example}{fg=white,bg=InkExample!88!white}
\setbeamercolor*{block body example}{fg=Ink,bg=InkExample!16!white}

% 页脚页码颜色（字号在各 style-*.tex 中设置）
\setbeamercolor{page number in head/foot}{fg=AcademicBlue}

% 文献表配色：作者与标题用黑色，出处与附注用灰色
\setbeamercolor{bibliography entry author}{fg=black}
\setbeamercolor{bibliography entry title}{fg=black}
\setbeamercolor{bibliography entry location}{fg=AcademicGray}
\setbeamercolor{bibliography entry note}{fg=AcademicGray}
     % 墨线配色与 Block 语义色
% =============================================================================
% layout.tex —— 版式层：页标题、边距、页脚与页码、过渡页、目录与列表模板
% -----------------------------------------------------------------------------
% 依赖：colors.tex（字体颜色）、style-*.tex（\InkMargin、\InkTopRule 等）。
% 本文件是「改版面观感」的主要入口。
% =============================================================================

% -----------------------------------------------------------------------------
% 1. 页标题与边距
% -----------------------------------------------------------------------------
% 标题、横线和正文共享同一组左右边距，只需在 style-*.tex 里改 \InkMargin。
\setbeamersize{text margin left=\InkMargin,text margin right=\InkMargin}

% 标题下方蓝色横线的粗细
\newcommand{\InkTopRule}{0.7pt}

% 自定义页标题模板：标题（+可选副标题）→ 蓝色细横线 → 下方留白。
% 用 \hrule 而不是 \hline，避免额外的表格上下文；
% 顶部 0.35cm 与底部 0.25cm 的留白决定正文起始位置。
\setbeamertemplate{frametitle}{%
  \nointerlineskip
  \begin{beamercolorbox}[wd=\textwidth,sep=0pt]{frametitle}
    \vskip0.35cm
    \usebeamerfont{frametitle}\strut\insertframetitle\strut\par
    \ifx\insertframesubtitle\empty\else
      {\usebeamerfont{framesubtitle}\insertframesubtitle\par}%
    \fi
    \vskip0.12cm
    {\color{AcademicBlue}\hrule height\InkTopRule}
    \vskip0.25cm
  \end{beamercolorbox}}

% 多页 frame 自动续页时，在标题后追加序号，例如「（2）」
\setbeamertemplate{frametitle continuation}{%
  \ifnum\insertcontinuationcount>1\relax\space(\insertcontinuationcount)\fi}

% 标题、章节页用衬线粗体（中文即思源宋体 Bold）
\setbeamerfont{title}{family=\rmfamily,series=\bfseries}
\setbeamerfont{frametitle}{family=\rmfamily,series=\bfseries}
\setbeamerfont{section title}{family=\rmfamily,series=\bfseries}

% 关闭右下角导航图标；文献表不显示条目图标
\setbeamertemplate{navigation symbols}{}
\setbeamertemplate{bibliography item}{}
% 图表编号
\setbeamertemplate{caption}[numbered]

% -----------------------------------------------------------------------------
% 2. 列表模板
% -----------------------------------------------------------------------------
\setbeamertemplate{enumerate item}{\insertenumlabel.}
\setbeamertemplate{enumerate subitem}{\insertenumlabel.\insertsubenumlabel.}
\setbeamertemplate{itemize item}{\small\textbullet}
\setbeamertemplate{itemize subitem}{\textendash}

% -----------------------------------------------------------------------------
% 3. 页脚与页码（含附录独立编号）
% -----------------------------------------------------------------------------
% 正文页码：「当前 / 总」。沿用 Beamer 内置的 totalframenumber 模板，
% 其内容等价于原来的 \insertframenumber\,/\,\inserttotalframenumber。
%
% 附录页码：独立编号「A-序号 / A-总数」。
%   \insertframenumberinappendix   附录内第几页（从 1 开始）
%   \insertappendixframenumber     附录总页数
% 这两个计数器由 Beamer 在 \appendix 时维护，并借助 .nav 辅助文件跨轮传递，
% 因此不需要手工硬编码任何数字，也不需要自己加计数器。
% 注意：绝不能在页脚里判断 \ifbeamer@inappendix —— 该条件在导言区
%       渲染页脚时还未就绪，会直接报 Undefined control sequence。
%       正确做法是用下面的 \appendix 包装，在进入附录时切换模板。
\defbeamertemplate*{page number in head/foot}{inkappendixframenumber}{%
  A-\insertframenumberinappendix\,/\,A-\insertappendixframenumber}

% 当前使用的页码模板名；\appendix 会把它切换成附录模板。
% 使用 \gdef 是因为 \appendix 在正文中执行，需要全局生效。
\def\InkPageNumberStyle{totalframenumber}

% 统一的页码输出：右对齐 + 当前模板。字号/颜色由 \ApplySlideStyle 与 colors.tex 设定。
% \setbeamertemplate 在这里位于每个页脚的分组内部，作用范围仅限本次排版；
% 由于模板名是全局的，切换一次即对所有后续页面稳定生效。
\newcommand{\InkPageNumber}{%
  \setbeamertemplate{page number in head/foot}[\InkPageNumberStyle]%
  \usebeamerfont{page number in head/foot}%
  \usebeamercolor[fg]{page number in head/foot}\hfill
  \usebeamertemplate***{page number in head/foot}\par}

% 包装 \appendix：先执行 Beamer 原本的附录处理（进入附录模式、写入计数器），
% 再把页码模板切换到附录版本。这样附录页脚显示 A-1 / A-3 式编号，
% 而正文始终保持「当前 / 总」不变。
\let\InkAppendixOriginal\appendix
\renewcommand{\appendix}{%
  \InkAppendixOriginal
  \global\def\InkPageNumberStyle{inkappendixframenumber}}

% -----------------------------------------------------------------------------
% 4. 目录
% -----------------------------------------------------------------------------
% 一级目录带编号（1. 研究背景）；二级目录用紧凑自定义模板：左缩进 2em，
% 形如「3.1 小节名」，避免默认模板占用的额外竖直空间。
\setbeamertemplate{section in toc}[sections numbered]
\setbeamertemplate{subsection in toc}{%
  \leavevmode\leftskip=2em\relax
  \inserttocsectionnumber.\inserttocsubsectionnumber\enspace
  \inserttocsubsection\par}

% -----------------------------------------------------------------------------
% 5. 章节与小节过渡页
% -----------------------------------------------------------------------------
% 写在 frame 之外的 \section / \subsection 会自动插入一张居中过渡页，
% 无需手写额外 frame。两类过渡页都不显示页码（noframenumbering），
% 因此不会占用正文页码，也不受附录编号影响。
% 带星号的 \section* / \subsection* 不触发过渡页。

% 章节页：水平垂直居中，显示「编号 + 标题」，例如「1. 研究背景」
\AtBeginSection[]{%
  \begin{frame}[plain,c,noframenumbering]
    \centering
    {\usebeamerfont{section title}\usebeamercolor[fg]{section title}%
      \insertsectionnumber.\enspace\insertsection\par}
  \end{frame}
}

% 小节页：上下结构。上方灰色「章节号 + 章节名」，下方较大的「章节号.小节号 + 小节名」
\AtBeginSubsection[]{%
  \begin{frame}[plain,c,noframenumbering]
    \centering
    {\normalsize\color{AcademicGray}%
      \insertsectionnumber.\enspace\insertsection\par}
    \vspace{0.8em}
    {\usebeamerfont{section title}\usebeamercolor[fg]{section title}\insertsectionnumber.\insertsubsectionnumber\enspace\insertsubsection\par}
  \end{frame}
}
     % 页标题、边距、页脚与页码、过渡页
% 目录：悬挂缩进、按实际高度分栏；双栏仍超高时交给 frame 自动续页。
% 依赖 layout.tex；在带 allowframebreaks 的目录 frame 中调用 \InkAutoToc。

% 编号占用固定宽度，小节与章节标题的文字起点对齐。
\newlength{\InkTocIndent}
\setlength{\InkTocIndent}{2.4em}
\newcommand{\InkTocNumber}{%
  {\setlength{\fboxsep}{0.22em}%
   \colorbox{AcademicBlue!9!white}{%
     \makebox[1.1em][c]{\sffamily\bfseries\color{AcademicBlue}%
       \strut\ifnum\inserttocsectionnumber<10 0\fi\inserttocsectionnumber}}}}

\setbeamertemplate{section in toc}{%
  \begingroup
  \par\vskip0.45em
  % 编号背景不参与行高计算，避免不同标题触发不同的基线间距。
  \setbox0=\hbox to\InkTocIndent{\smash{\InkTocNumber}\hfil}%
  \hangindent=\wd0\hangafter=1
  \leavevmode\box0\strut\inserttocsection\par
  \endgroup}
\setbeamertemplate{subsection in toc}{%
  \begingroup
  \par\vskip0.3em\nointerlineskip
  \leftskip=\InkTocIndent
  \setbox0=\hbox{{\color{AcademicGray}\inserttocsectionnumber.\inserttocsubsectionnumber}\enspace}%
  \hangindent=\wd0\hangafter=1
  \leavevmode\box0\strut\inserttocsubsection\par
  \endgroup}

\newsavebox{\InkTocBox}
\newsavebox{\InkTocTitleBox}
\newsavebox{\InkTocLeft}
\newsavebox{\InkTocRight}
\newsavebox{\InkTocBestLeft}
\newsavebox{\InkTocBestRight}
\newdimen\InkTocAvailable
\newdimen\InkTocHeight
\newdimen\InkTocBestHeight
\newcount\InkTocLastSection
\newcount\InkTocSplit
\newcount\InkTocRightStart

% #1 盒子寄存器，#2 排版宽度，#3 tableofcontents 的选项。
\newcommand{\InkMeasureToc}[3]{%
  \setbox#1=\vbox{%
    \hsize=#2\relax\linewidth=\hsize
    % 抵消 Beamer 目录自带的负顶部间距。
    \vskip0.5em\tableofcontents[#3]}}

\newcommand{\InkAutoToc}{%
  \begingroup
  \setlength{\parskip}{0pt}%
  \global\InkTocLastSection=0
  \addtobeamertemplate{section in toc}{%
    \ifnum\inserttocsectionnumber>\InkTocLastSection
      \global\InkTocLastSection=\inserttocsectionnumber\relax
    \fi}{}%
  \InkMeasureToc{\InkTocBox}{\textwidth}{}%
  \setbox\InkTocTitleBox=\vbox{\hsize=\textwidth\usebeamertemplate{frametitle}}%
  % allowframebreaks 默认每页使用 95% 的可用高度，预留相同余量。
  \InkTocAvailable=\dimexpr\textheight-\ht\InkTocTitleBox-\dp\InkTocTitleBox-0.6cm\relax
  \InkTocAvailable=0.95\InkTocAvailable
  \par\vspace*{0.15cm}%
  \ifdim\dimexpr\ht\InkTocBox+\dp\InkTocBox\relax>\InkTocAvailable
    \InkTocBestHeight=\maxdimen
    \InkTocSplit=1
    % 只在章节边界拆分，选取较高一栏最矮的位置，避免小节分布不均。
    \loop\ifnum\InkTocSplit<\InkTocLastSection
      \InkTocRightStart=\InkTocSplit \advance\InkTocRightStart by1
      \InkMeasureToc{\InkTocLeft}{0.48\textwidth}{sections={1-\the\InkTocSplit}}%
      \InkMeasureToc{\InkTocRight}{0.48\textwidth}{sections={\the\InkTocRightStart-\the\InkTocLastSection}}%
      \InkTocHeight=\dimexpr\ht\InkTocLeft+\dp\InkTocLeft\relax
      \ifdim\dimexpr\ht\InkTocRight+\dp\InkTocRight\relax>\InkTocHeight
        \InkTocHeight=\dimexpr\ht\InkTocRight+\dp\InkTocRight\relax
      \fi
      \ifdim\InkTocHeight<\InkTocBestHeight
        \InkTocBestHeight=\InkTocHeight
        \setbox\InkTocBestLeft=\copy\InkTocLeft
        \setbox\InkTocBestRight=\copy\InkTocRight
      \fi
      \advance\InkTocSplit by1
    \repeat
    \ifdim\InkTocBestHeight>\InkTocAvailable
      % 不缩小字号或隐藏条目；释放单栏盒子，允许在条目之间分页。
      \unvbox\InkTocBox
    \else
      \noindent
      \begin{minipage}[t]{0.48\textwidth}
        \vspace{0pt}\unvbox\InkTocBestLeft
      \end{minipage}\hfill
      \begin{minipage}[t]{0.48\textwidth}
        \vspace{0pt}\unvbox\InkTocBestRight
      \end{minipage}\par
    \fi
  \else
    \noindent\usebox{\InkTocBox}\par
  \fi
  \endgroup}
        % 目录换行、自适应分栏与续页
% =============================================================================
% blocks.tex —— 三色圆角 Block（普通 / 警示 / 示例）
% -----------------------------------------------------------------------------
% 目标：给 Beamer 原生圆角盒加上内边距，并让块内的列表符号与 \alert
%       跟随该块自己的颜色（而不是全局主题蓝）。
% 做法：不引入任何新宏包，只用 beamerboxesrounded + minipage 内缩。
% 依赖：colors.tex 必须先加载（用到 block title/body 三套配色）。
% 调参入口：\InkBlockPadding 控制水平内边距；两个 \vspace 控制竖直内边距。
% =============================================================================

% 标题栏与正文的左右内边距
\newlength{\InkBlockPadding}
\setlength{\InkBlockPadding}{0.5em}

% 内层 minipage 的可用宽度 = 行宽 - 左右内边距（只需算一次，避免重复算式）
\newlength{\InkBlockInnerWidth}

% 把 block 内的「局部结构色」和「强调色」重绑到当前块的标题色，
% 于是块内的 itemize/enumerate 符号和 \alert 会自动变成该块的颜色。
% 参数：空 / " alerted" / " example"（注意前导空格，用于拼接模板名）
\newcommand{\InkBlockColors}[1]{%
  \setlength{\parskip}{0.3em}%
  \setbeamercolor{local structure}{use={block title#1},fg=block title#1.bg}%
  \setbeamercolor{alerted text}{use={block title#1},fg=block title#1.bg}}

% 块开始：先排一个自带内边距的标题栏，再开启正文 minipage。
% shadow=false 与规范所述「无阴影」一致；width=\linewidth 让块自适应单栏或分栏。
\newcommand{\InkBlockBegin}[1]{%
  \par\vskip\medskipamount
  % 原生圆角盒在左右各向外延伸 4bp；内缩后让外框与正文边界对齐。
  \noindent\hspace*{4bp}%
  \begin{minipage}{\dimexpr\linewidth-8bp\relax}
  \setlength{\InkBlockInnerWidth}{\dimexpr\linewidth-\InkBlockPadding-\InkBlockPadding\relax}%
  \begin{beamerboxesrounded}[upper=block title#1,lower=block body#1,width=\linewidth,shadow=false]{%
    \noindent\hspace*{\InkBlockPadding}%
    \begin{minipage}{\InkBlockInnerWidth}
      \usebeamerfont{block title#1}\vspace*{0.2ex}%
      \strut\insertblocktitle\par\vspace*{0.2ex}%
    \end{minipage}}
    \usebeamerfont{block body#1}\InkBlockColors{#1}%
    \noindent\hspace*{\InkBlockPadding}%
    \begin{minipage}{\InkBlockInnerWidth}
      \vspace*{0.6ex}}

% 块结束：补上正文下边距并闭合两层盒子
\newcommand{\InkBlockEnd}{%
      \par\vspace*{0.6ex}%
    \end{minipage}\end{beamerboxesrounded}%
  \end{minipage}\par\vskip\smallskipamount}

% 把上面两个宏挂到 Beamer 的 six 个 Block 模板上
\setbeamertemplate{block begin}{\InkBlockBegin{}}
\setbeamertemplate{block end}{\InkBlockEnd}
\setbeamertemplate{block alerted begin}{\InkBlockBegin{ alerted}}
\setbeamertemplate{block alerted end}{\InkBlockEnd}
\setbeamertemplate{block example begin}{\InkBlockBegin{ example}}
\setbeamertemplate{block example end}{\InkBlockEnd}

% 块内文字：无衬线（黑体）粗体标题 + 常规正文；
% alerted / example 的字体继承普通块，只换配色。
\setbeamerfont{block title}{family=\sffamily,series=\bfseries}
\setbeamerfont{block body}{family=\sffamily,size=\normalsize,series=\mdseries,shape=\upshape}
\setbeamerfont{block title alerted}{parent=block title}
\setbeamerfont{block title example}{parent=block title}
\setbeamerfont{block body alerted}{parent=block body}
\setbeamerfont{block body example}{parent=block body}
     % 三色圆角 Block
% =============================================================================
% commands.tex —— 目录、脚注等「分散在各处的小设置」
% -----------------------------------------------------------------------------
% 依赖：base.tex（\SourceNoteSize 等）、colors.tex、layout.tex。
% 页码相关的 \InkPageNumber 定义在 layout.tex（与页脚模板放在一起）。
% =============================================================================

% -----------------------------------------------------------------------------
% 1. 目录（两级）字体
% -----------------------------------------------------------------------------
% 一级目录：正文尺寸 + 衬线 + Slanted；二级目录：小字号 + 衬线 + Slanted。
% 具体字号在各 style-*.tex 中覆盖（小版 \scriptsize，大版 \footnotesize），
% 这里只固定字族与字形，保证两版一致。
\setbeamerfont{section in toc}{size=\normalsize,family=\rmfamily,shape=\slshape}
\setbeamerfont{subsection in toc}{family=\rmfamily,shape=\slshape}

% -----------------------------------------------------------------------------
% 2. 脚注
% -----------------------------------------------------------------------------
% 脚注正文：小字号、衬线、Italic；脚注编号保持直立（斜体数字难辨认）。
% 字号由各 style-*.tex 设置（小版 \scriptsize，大版 \footnotesize）。
\setbeamerfont{footnote}{family=\rmfamily,shape=\itshape}
\setbeamerfont{footnote mark}{shape=\upshape}
   % 目录与脚注字体
\ApplySlideStyle          % 本版式的字号与页脚
\title{调研主题或项目名称}
\date{\today}
   % 标题与日期（封面字段）

% 正文段间距；目录页内单独归零。
\setlength{\parskip}{0.5em}
\begin{document}

% -----------------------------------------------------------------------------
% 3. 封面
%    局部分组加一段脚注后的固定留白：Beamer 把脚注排在底部，
%    plain 封面容易贴边，把脚注文字上移 0.4cm。分组保证正文页脚注不受影响。
% -----------------------------------------------------------------------------
\begingroup
\addtobeamertemplate{footnote}{}{\par\vspace*{0.4cm}}
\begin{frame}[plain,c,noframenumbering]
  \titlepage
\end{frame}
\endgroup

% -----------------------------------------------------------------------------
% 4. 目录（不计入页码）
% -----------------------------------------------------------------------------
\begin{frame}[plain,allowframebreaks,noframenumbering]{目录}
  \InkAutoToc
\end{frame}

% -----------------------------------------------------------------------------
% 5. 语法示例手册：按功能查找，复制所需 frame 即可用于自己的汇报。
% -----------------------------------------------------------------------------
% 示例阅读顺序：查看实际效果，在源码中复制对应 frame，再按需替换内容。
\section{文字与列表}
\subsection{页面与文字}

% 含 verbatim 的 frame 必须使用 fragile；示例中的结束行带注释，避免被外层提前识别。
\begin{frame}[fragile]{页面与章节的基本写法}
  \begin{verbatim}
\section{章节名}
\subsection{小节名}
\begin{frame}{页面标题}
  页面正文。
\end{frame} % 结束示例页面
  \end{verbatim}
  \slideNote{章节和小节会自动生成过渡页。星号形式不进入目录，
  也不触发过渡页；封面标题与日期在 config/metadata.tex 中设置。}
\end{frame}

\begin{frame}{文字、强调与页面副标题}
  \framesubtitle{副标题使用 framesubtitle 命令设置}
  正文支持中文与 English 混排，段落之间使用空行分隔。

  \textbf{粗体}用于强调；\textit{Italic text}展示西文斜体；
  \texttt{inline code}用于代码片段。

  \keypoint{主题色粗体重点}，以及 \alert{需要强调的条件}。
  \slideNote{正文保留黑体；页标题和表格使用宋体。注释文字自动使用小一号字号，
  与正文保持清晰的层级。}
\end{frame}

\subsection{列表与代码}
% 嵌套列表：在 item 内再放入 itemize / enumerate 环境。
\begin{frame}{无序、有序与嵌套列表}
  \begin{columns}[T,onlytextwidth]
    \column{0.48\textwidth}
      \textbf{无序列表：itemize}
      \begin{itemize}
        \item 第一项。
          \begin{itemize}
            \item 嵌套子项。
            \item 另一条子项。
          \end{itemize}
        \item 第二项。
      \end{itemize}
    \column{0.48\textwidth}
      \textbf{有序列表：enumerate}
      \begin{enumerate}
        \item 第一步。
          \begin{enumerate}
            \item 第一步的子步骤。
          \end{enumerate}
        \item 第二步。
      \end{enumerate}
  \end{columns}
  \slideNote{使用 item 添加条目；在列表内部再放一个列表环境即可嵌套。}
\end{frame}

\begin{frame}{宽松要点与术语列表}
  \begin{columns}[T,onlytextwidth]
    \column{0.48\textwidth}
      \textbf{keypoints：宽松要点}
      \begin{keypoints}
        \item 第一个核心观点。
        \item 第二个核心观点。
        \item 第三个核心观点。
      \end{keypoints}
    \column{0.48\textwidth}
      \textbf{description：术语说明}
      \begin{description}[模型]
        \item[数据] 输入与观测。
        \item[模型] 结构与参数。
        \item[指标] 评价的依据。
      \end{description}
  \end{columns}
\end{frame}

\begin{frame}{网址、路径与特殊字符}
  \begin{itemize}
    \item 网址：\url{https://ctan.org/pkg/beamer}
    \item 命名链接：\href{https://ctan.org/pkg/beamer}{Beamer 文档入口}
    \item 路径：\path{figures/example-figure.png}
    \item 行内代码：\texttt{tectonic main.tex}
  \end{itemize}

  \begin{block}{特殊字符需要转义}
    百分号写 \texttt{\textbackslash\%}，与号写 \texttt{\textbackslash\&}，
    下划线写 \texttt{\textbackslash\_}，井号写 \texttt{\textbackslash\#}。
    需要保留原始排版的整段代码，可把 frame 声明为
    \texttt{\textbackslash begin\{frame\}[fragile]} 后使用 verbatim 环境。
  \end{block}
\end{frame}

\begin{frame}[fragile]{整段代码与原样文本}
  页面声明：\verb|\begin{frame}[fragile]{代码示例}|。
  \begin{verbatim}
score = mean(values)
# keep spaces and symbols: _ % &
  \end{verbatim}
  \slideNote{verbatim 保留空格和特殊字符；含有原样代码的页面必须使用 fragile。
  不需要额外的语法高亮宏包或外部程序。}
\end{frame}

% 示例阅读顺序：查看实际效果，在源码中复制对应 frame，再按需替换内容。
\section{Block 与注释}
\subsection{三色块与组合}
\begin{frame}{Block 样式示例}
  \begin{block}{普通说明：block}
    用于研究问题、定义或方法说明。正文可以包含\textbf{粗体}与公式，例如 $E=mc^2$。
  \end{block}

  \begin{alertblock}{注意事项：alertblock}
    用于假设、限制和需要关注的结论；可用 \alert{强调文字}标出关键条件。
  \end{alertblock}

  \begin{exampleblock}{具体示例：exampleblock}
    用于算例、实验设置或应用场景。例如：对三组数据采用相同指标进行比较。
  \end{exampleblock}
\end{frame}

% 三色环境：block / alertblock / exampleblock；分栏使用 [T,onlytextwidth]。
\begin{frame}{Block 中的列表与双栏}
  \begin{columns}[T,onlytextwidth]
    \column{0.48\textwidth}
    \begin{block}{方法步骤}
      \begin{enumerate}
        \item 明确研究对象。
        \item 统一评价指标。
        \item 对比实验结果。
      \end{enumerate}
    \end{block}

    \column{0.48\textwidth}
    \begin{alertblock}{比较条件}
      \begin{itemize}
        \item 使用相同数据集。
        \item 记录参数设置。
        \item 报告适用范围。
      \end{itemize}
    \end{alertblock}
  \end{columns}

  \begin{exampleblock}{带公式的示例}
    样本均值为 $\bar{x}=\frac{1}{n}\sum_{i=1}^{n}x_i$。
    对样本 $1,2,3$，有 $\bar{x}=2$。
  \end{exampleblock}
\end{frame}

% 空字符串和纯空格均视为无标题；三类环境保留各自的正文背景色。
\begin{frame}{无标题 Block}
  \begin{block}{}
    普通信息：空标题不产生标题栏，直接显示正文。
  \end{block}
  \begin{alertblock}{   }
    注意条件：纯空格标题也不会产生空白标题栏。
  \end{alertblock}
  \begin{exampleblock}{}
    应用示例：支持 $x>0$、\textbf{粗体}和多行文字。
  \end{exampleblock}
\end{frame}

\subsection{独立段落注释}
% 用法：\inknote{补充说明}；结束后恢复正文样式。
\begin{frame}[fragile]{正文中的段落注释}
  这是正常正文，用于陈述主要观点和分析结果。
  \inknote{这是一段补充注释：可以解释假设、范围或阅读提示。
  它没有标题栏，采用深蓝细竖线、淡灰直角背景与小字号深灰衬线文字，与正文自然区分。}
  注释结束后，正文恢复原来的字号、颜色和字重。
  \begin{verbatim}
\inknote{这里填写注释内容。}
  \end{verbatim}
\end{frame}

\begin{frame}{分栏注释与来源说明}
  \begin{columns}[T,onlytextwidth]
    \column{0.48\textwidth}
      \textbf{方法说明}\par\smallskip
      这里放置正文段落。
      \inknote{注释自动适应当前栏宽，也可包含 $n>0$、\textbf{粗体}和链接。

      空行可分隔注释内部的段落。}
    \column{0.48\textwidth}
      \textbf{结果说明}\par\smallskip
      这里放置另一段正文。
      \inknote{inknote 用于解释正文；inksource 用于简短的图表出处。
      两者的语义和视觉层级不同。}
  \end{columns}
  \inksource{来源：此处填写数据来源或参考文献。}
\end{frame}

% 示例阅读顺序：查看实际效果，在源码中复制对应 frame，再按需替换内容。
\section{分栏与图片}
\subsection{分栏排版}
\begin{frame}{双栏与不等宽分栏}
  \begin{columns}[T,onlytextwidth]
    \column{0.56\textwidth}
      \textbf{较宽的文字栏}
      \begin{enumerate}
        \item 使用 columns 创建分栏。
        \item 用 column 指定每栏宽度。
        \item 两栏之间预留适当空隙。
      \end{enumerate}
    \column{0.40\textwidth}
      \begin{block}{较窄的内容栏}
        放置结论、公式或图片，顶部与左栏对齐。
      \end{block}
  \end{columns}
  \slideNote{使用 [T,onlytextwidth] 顶部对齐并限制总宽度；
  栏宽之和小于正文宽度，为栏间留出空隙。}
\end{frame}

\subsection{插图与图题}
% 用法：\slideimage[高度上限]{文件名}；来源说明放在对应图片所在栏内。
\begin{frame}{插图与来源标注}
  \begin{columns}[T,onlytextwidth]
    \column{0.54\textwidth}
      \centering
      \slideimage[0.5\textheight]{example-figure.png}
      \sourceNote{来源：模板内置占位图。}
    \column{0.42\textwidth}
      \begin{itemize}
        \item \texttt{\textbackslash slideimage\{文件名\}}：按栏宽等比缩放。
        \item 加可选参数可限制高度，例如 \texttt{[0.45\textbackslash textheight]}。
        \item 图片统一放在 \texttt{figures/}，只写文件名即可。
      \end{itemize}
  \end{columns}
\end{frame}

\begin{frame}{并排两图}
  \begin{columns}[T,onlytextwidth]
    \column{0.48\textwidth}
      \centering
      \begin{figure}
        \slideimage[0.36\textheight]{example-figure.png}
        \caption{方案 A}
      \end{figure}
    \column{0.48\textwidth}
      \centering
      \begin{figure}
        \slideimage[0.36\textheight]{example-figure.png}
        \caption{方案 B}
      \end{figure}
  \end{columns}
  \sourceNote{来源：同一占位图重复使用，正式汇报时替换为两组真实结果。}
\end{frame}

% caption 之后添加 label；通过 ref 引用实际编号，不手写图号。
\begin{frame}{带图题与交叉引用的图片}
  \begin{figure}
    \centering
    \slideimage[0.36\textheight]{example-figure.png}
    \caption{示例图片与图题}\label{fig:demo}
    \sourceNote{来源：模板内置占位图。}
  \end{figure}
  图~\ref{fig:demo} 使用 figure、caption 和 label 添加图题与编号引用。
\end{frame}

\section{表格}
\subsection{三线表与表头}
\begin{frame}[fragile]{三线表与粗体表头}
  \begin{tabular}{@{}lrr@{}}
    \toprule
    \tablehead{方法} & \tablehead{准确率} & \tablehead{耗时（秒）} \\
    \midrule
    基线 & 0.81 & 12.3 \\
    改进 & 0.87 & 10.8 \\
    \bottomrule
  \end{tabular}
  \begin{verbatim}
\tablehead{方法} & \tablehead{准确率} \\
  \end{verbatim}
  \slideNote{表头每个单元格用 tablehead 标记，统一使用衬线粗体。
  l、c、r 分别控制左、中、右对齐；三线表使用 top/mid/bottomrule。}
\end{frame}

\subsection{自适应宽度与合并表头}
\begin{frame}{自动列宽与多行内容}
  \small
  \begin{tabularx}{\linewidth}{@{}lYY@{}}
    \toprule
    \tablehead{方法} & \tablehead{适用场景} & \tablehead{注意事项} \\
    \midrule
    文献调研 & 归纳不同工作的共同问题、方法与发展趋势 & 统一比较口径，注明资料的时间范围 \\
    系统调研 & 检查公开实现与实验设置，进行复现 & 记录版本、运行环境和关键参数 \\
    \bottomrule
  \end{tabularx}
  \sourceNote{示例内容；Y 列自动分配宽度并换行，文字左对齐。}
  \slideNote{表格行距统一为 1.2。使用 tabularx 与 Y 列处理长文字，
  不需要手工在每个单元格中插入换行。}
\end{frame}

\begin{frame}{合并表头与表格引用}
  \begin{table}
    \centering
    \small
    \begin{tabular}{@{}lrrrr@{}}
      \toprule
      & \multicolumn{2}{c}{\tablehead{数据集 A}} & \multicolumn{2}{c}{\tablehead{数据集 B}} \\
      \cmidrule(lr){2-3}\cmidrule(l){4-5}
      \tablehead{方法} & \tablehead{均值} & \tablehead{标准差} & \tablehead{均值} & \tablehead{标准差} \\
      \midrule
      基线 & 0.81 & 0.03 & 0.79 & 0.04 \\
      改进 & 0.87 & 0.02 & 0.85 & 0.03 \\
      \bottomrule
    \end{tabular}
    \caption{两组数据集上的示例结果}\label{tab:demo}
  \end{table}
  表~\ref{tab:demo} 演示 multicolumn 合并单元格与 cmidrule 分组横线。
  \sourceNote{数值为排版示例，不代表实际实验结果。}
\end{frame}
\begin{frame}{分栏与表格}
  \begin{columns}[T,onlytextwidth]
    \column{0.50\textwidth}
      \textbf{要点}（左栏）
      \begin{enumerate}
        \item 栏宽按内容分配，两栏合计不超过 $1$。
        \item \texttt{[T]} 让两栏顶部对齐；\texttt{onlytextwidth} 避免整体溢出。
        \item 栏内可放列表、公式或图片。
      \end{enumerate}
    \column{0.46\textwidth}
      \small
      \begin{tabularx}{\linewidth}{@{}lY@{}}
        \toprule
        \tablehead{维度} & \tablehead{说明} \\
        \midrule
        数据 & 填写数据来源与规模 \\
        指标 & 填写评价指标与口径 \\
        结论 & 填写该维度的结论 \\
        \bottomrule
      \end{tabularx}
      \sourceNote{表格使用 \texttt{tabularx} 与 \texttt{Y} 列自适应宽度。}
  \end{columns}
\end{frame}

% 示例阅读顺序：查看实际效果，在源码中复制对应 frame，再按需替换内容。
\section{公式与分步演示}
\subsection{公式与数学环境}
\begin{frame}{行内与行间公式}
  \begin{block}{公式写法}
    行内公式写作 $a^2+b^2=c^2$；需要居中并强调时使用行间形式：
    \[
      \bar{x}=\frac{1}{n}\sum_{i=1}^{n}x_i,\qquad
      s^2=\frac{1}{n-1}\sum_{i=1}^{n}(x_i-\bar{x})^2 .
    \]
  \end{block}

  \begin{exampleblock}{常用记号}
    均值 $\bar{x}$、方差 $s^2$、求和 $\sum_{i=1}^{n}$、极限 $\lim_{n\to\infty}$。
  \end{exampleblock}
\end{frame}

% align 用 & 对齐，eqref 引用公式；pmatrix 与 cases 分别提供矩阵和分段函数。
\begin{frame}{多行公式、矩阵与分段函数}
  \begin{align}
    \label{eq:demo} f(x) &= (x+1)^2 \\
         &= x^2+2x+1.\notag
  \end{align}
  式~\eqref{eq:demo} 使用 align 对齐，label 与 eqref 用于引用。
  \[
    A=\begin{pmatrix}1&2\\3&4\end{pmatrix},\qquad
    g(x)=\begin{cases}x,&x\geq0,\\-x,&x<0.\end{cases}
  \]
  \inknote{pmatrix 提供带圆括号的矩阵，cases 用于分段函数。
  使用 \texttt{\&} 设置对齐点，双反斜线换行。}
\end{frame}

\begin{frame}{定义、定理与证明}
  \begin{definition}[样本均值]
    对 $n\geq1$ 个实数 $x_1,\ldots,x_n$，定义
    $\bar{x}=\frac{1}{n}\sum_{i=1}^{n}x_i$。
  \end{definition}

  \begin{theorem}[均值界]
    $\min_i x_i\leq\bar{x}\leq\max_i x_i$。
  \end{theorem}

  \begin{proof}
    对每个 $i$，均有 $\min_j x_j\leq x_i\leq\max_j x_j$。
    将不等式求和，再除以 $n$ 即得结论。
  \end{proof}
\end{frame}

\subsection{逐步展示与链接}
\begin{frame}{分步显示}
  \begin{exampleblock}{用 \texttt{\textbackslash pause} 逐条显示}
    \begin{enumerate}
      \item 第一步：先给出问题定义。
      \pause
      \item 第二步：再给出方法。
      \pause
      \item 第三步：最后给出结果与结论。
    \end{enumerate}
  \end{exampleblock}
\end{frame}

\begin{frame}{按步骤强调}
  \begin{itemize}
    \item 基线方法：按固定顺序处理输入。
    \item \alert<2>{改进点一}：在第一步加入筛选。
    \item \alert<3>{改进点二}：在第二步加入校准。
    \item \alert<4>{结论}：两项改进共同带来提升。
  \end{itemize}
\end{frame}

% uncover 保留空间；only 不保留隐藏内容；overlayarea 固定替换区域高度。
\begin{frame}[label=overlay-demo]{固定布局与替换内容}
  \begin{columns}[T,onlytextwidth]
    \column{0.48\textwidth}
      \textbf{uncover：保留空间}
      \begin{itemize}
        \item<1-> 第一步出现后保留。
        \item<2-> 第二步出现后保留。
      \end{itemize}
      \uncover<3->{第三步显示这段补充说明。}
    \column{0.48\textwidth}
      \textbf{only：替换内容}
      \begin{overlayarea}{\linewidth}{2.8em}
        \only<1>{第一步：准备输入。}
        \only<2>{第二步：计算结果。}
        \only<3>{第三步：分析输出。}
      \end{overlayarea}
  \end{columns}
  \inknote{uncover 保留隐藏文字的位置；only 只排版当前步骤的内容。
  overlayarea 固定替换区域的高度，避免页面跳动。}
  \hyperlink{overlay-demo<1>}{\beamerbutton{返回第一步}}
\end{frame}

% 示例阅读顺序：查看实际效果，在源码中复制对应 frame，再按需替换内容。
\section{引文与文献}
\subsection{引用文字与出处}
\begin{frame}{引文与引用}
  \begin{quote}
    \enquote{引文内容写在这里，通常是一句原文；\texttt{quote} 环境整段缩进，
    适合较短的引文。}
  \end{quote}

  \vspace{0.3em}
  \begin{itemize}
    \item 叙述式引用：\textcite{example2026} 提出了该问题的基本框架。
    \item 括号式引用：已有工作给出了类似结论\parencite{example2026}。
    \item 只取作者与年份：\citeauthor{example2026}（\citeyear{example2026}）。
    \item 长引文用 \texttt{\textbackslash blockquote}：\blockquote{这是一段
      较长的引文，超过一定行数后会自动转为整段缩进，并在末尾给出出处。}
  \end{itemize}
  \sourceNote{来源：\parencite{example2026}（示例条目，请替换为真实文献）。}
\end{frame}

\begin{frame}{多段引文与普通脚注}
  \begin{quotation}
    第一段引文：quotation 适合包含多个段落的引用文字。

    第二段引文：空行分段，保留引文环境的缩进和专用字体。
  \end{quotation}
  正文中的补充说明可以使用普通脚注\footnote{这是一条补充说明；脚注自动编号。}。
  \sourceNote{出处使用 sourceNote；需要文献时可加入 \parencite{example2026}。}
\end{frame}

\subsection{脚注引用与文献表}
% 普通脚注用 footnote，整条文献用 footcite；文献键来自 references.bib。
\begin{frame}{脚注式引用}
  \begin{itemize}
    \item 作者 + 年份：此处引用一条工作\footnote{\citeauthor{example2026},
      \citeyear{example2026}.}。
    \item 作者 + 标题 + 年份：此处引用另一条工作\footnote{\citeauthor{example2026},
      \citetitle{example2026}, \citeyear{example2026}.}。
    \item 完整文献条目：用 \texttt{\textbackslash footcite} 直接把文献表内容
      放进脚注\footcite{example2026}。
  \end{itemize}

  \begin{alertblock}{使用建议}
    \begin{itemize}
      \item 图表与数据来源用 \texttt{\textbackslash sourceNote}；
      \item 观点出处用 \texttt{\textbackslash parencite} / \texttt{\textbackslash textcite}；
      \item 只在需要展开出处细节（如封面引文）时使用 \texttt{\textbackslash footnote}。
    \end{itemize}
  \end{alertblock}
\end{frame}

\begin{frame}[fragile]{文献表与附录的写法}
  \begin{verbatim}
\appendix
\begin{frame}[allowframebreaks]{参考文献}
  \printbibliography[heading=none]
\end{frame} % 结束参考文献页面
  \end{verbatim}
  \slideNote{在 references/references.bib 中添加条目，正文通过键名引用。
  附录页码独立编号；实际的参考文献表在文档末尾展示。}
\end{frame}


% -----------------------------------------------------------------------------
% 6. 附录（可选）
%    \appendix 之后页码切换为 A-当前 / A-总数（见 config/layout.tex）。
%    参考文献与补充材料都在 90-appendix.tex；不需要附录时删除下面两行。
% -----------------------------------------------------------------------------
\appendix
% 附录：\appendix 由 main.tex 发出，附录页码独立计数，
% 页码切换为 A-当前 / A-总数（见 inkline/layout.tex）。
% 不需要附录时，把 main.tex 里的 \appendix 与该文件的 \input 一起删除即可。
\section{附录}

\begin{frame}[allowframebreaks]{参考文献}
  % 文献表字体由 inkline/base.tex 的 \bibfont 统一设为 Times 风格衬线体，
  % 中文条目仍由 ctex 回退到思源宋体，无需在此另行设置。
  \small
  \printbibliography[heading=none]
\end{frame}

% 术语与符号表：首次出现的缩写或记号可在此统一解释，避免正文反复展开
\begin{frame}{术语与符号}
  \small
  \begin{tabularx}{\linewidth}{@{}lY@{}}
    \toprule
    \textbf{符号 / 缩写} & \textbf{含义} \\
    \midrule
    符号 A & 在此给出定义与单位 \\
    缩写 B & 全称及首次出现位置 \\
    缩写 C & 全称及首次出现位置 \\
    \bottomrule
  \end{tabularx}
  \inksource{符号表按正文出现顺序排列；正式使用时替换为实际记号。}
\end{frame}

% framebreak 手动分页；内容过长时 allowframebreaks 也可自动分页。
\begin{frame}[allowframebreaks]{长内容续页}
  \textbf{第一页}
  \begin{itemize}
    \item allowframebreaks 允许一个 frame 分成多页。
    \item framebreak 可在需要的位置手动分页。
  \end{itemize}
  \inknote{续页标题会追加序号。此示例放在附录，以便同时查看附录页码。}
  \framebreak
  \textbf{第二页}
  \begin{itemize}
    \item 长文献列表可按相同方式自动分页。
    \item 分步动画请使用普通 frame，避免与自动分页混用。
  \end{itemize}
\end{frame}

\begin{frame}[plain,c,noframenumbering]
  \centering
  {\usebeamerfont{title}\usebeamercolor[fg]{title}示例结束\par}
  \vspace{0.8em}
  普通汇报可在这里放置结束语。
  \inknote{plain 隐藏页眉页脚；c 让内容垂直居中；
  noframenumbering 使此页不占用页码。}
\end{frame}


\end{document}
